% 工尺谱示例 — 孟浩然《春晓》配工尺谱
% 演示 luatex-cn 的 \工尺 / \音 / \拖音 命令
%
% API 速查：
%   \工尺{主字}{右侧工尺串}                    主字 + 工尺旁注
%   \工尺[缩放=0.6, 颜色=red]{主}{尺工}        参数化
%   \音[板]{尺}                                 工尺字 + 红色板点
%   \音[眼]{尺}                                 工尺字 + 红色眼点
%   \音[高]{尺} \音[低]{尺}                     八度前缀 (亻/彳)
%   \音[板, 高]{尺}                             板眼+八度双槽
%   \拖音[眼]                                   续拍单元 (延音线)
%   \工尺设置{颜色=black, 缩放=0.55}             全局默认
%
\documentclass{ltc-guji}

\title{春\空格 晓}

\chapter{孟浩然·工尺谱}

\begin{document}
\begin{正文}

% 全局默认：工尺字略大、黑色
\工尺设置{缩放=0.55}

% ------------------------------------------------------------------
% 第一句 — 春眠不觉晓
%   春(尺,2 拍) 眠(工,1) 不(六,1) 觉(五,1) 晓(尺,3)
% ------------------------------------------------------------------
\工尺{春}{\音[板]{尺}\拖音[眼]}%
\工尺{眠}{\音[板]{工}}%
\工尺{不}{\音[眼]{六}}%
\工尺{觉}{\音[板]{五}}%
\工尺{晓}{\音[眼]{尺}\拖音[板]\拖音[眼]}

% ------------------------------------------------------------------
% 第二句 — 处处闻啼鸟
%   处(六) 处(五) 闻(工) 啼(尺,2) 鸟(六,高八度)
% ------------------------------------------------------------------
\工尺{处}{\音[板]{六}}%
\工尺{处}{\音[眼]{五}}%
\工尺{闻}{\音[板]{工}}%
\工尺{啼}{\音[眼]{尺}\拖音[板]}%
\工尺{鸟}{\音[眼, 高]{六}}

% ------------------------------------------------------------------
% 第三句 — 夜来风雨声
%   夜(尺) 来(工) 风(六,2) 雨(五) 声(尺,3)
% ------------------------------------------------------------------
\工尺{夜}{\音[板]{尺}}%
\工尺{来}{\音[眼]{工}}%
\工尺{风}{\音[板]{六}\拖音[眼]}%
\工尺{雨}{\音[板]{五}}%
\工尺{声}{\音[眼]{尺}\拖音[板]\拖音[眼]}

% ------------------------------------------------------------------
% 第四句 — 花落知多少
%   花(工) 落(六,低八度) 知(五) 多(尺) 少(尺,4 拍·终止)
% ------------------------------------------------------------------
\工尺{花}{\音[板]{工}}%
\工尺{落}{\音[眼, 低]{六}}%
\工尺{知}{\音[板]{五}}%
\工尺{多}{\音[眼]{尺}}%
\工尺{少}{\音[板]{尺}\拖音[眼]\拖音[板]\拖音[眼]}

\end{正文}
\end{document}
